\exercise{12.9}{A positive lower bound and products tending upward}
\goal Prove the product assertions in parts (a) and (b), and identify
Exercise~12.7 as the special case requested in part (c).

\step{A preliminary fact: extracting one positive lower bound.}
Write $\alpha=\liminf t_n$, and suppose $\alpha>0$. The value
$\alpha=+\infty$ is allowed. We claim that there are a real number
$c>0$ and an integer $N_0$ such that
\[
 t_n\ge c\qquad(n>N_0).
\]
If $\alpha$ is finite, choose $c=\alpha/2$; then $0<c<\alpha$.
If $\alpha=+\infty$, choose $c=1$ instead. In either case, the tail
infima $\inf\{t_n:n>N\}$ eventually exceed $c$, by their convergence
to $\alpha$. Each term in such a tail is at least its infimum, and
therefore exceeds $c$. Replacing this strict inequality by $\ge c$
gives the claim. This argument also covers an initially unbounded
sequence: since $\liminf t_n>0$, some tail has a finite lower bound.

\step{(a) An ordinary limit of $+\infty$.}
Assume $s_n\to+\infty$. To prove $s_nt_n\to+\infty$, fix an
arbitrary real threshold $M>0$. By the definition of this limit,
there is $N_1$ such that
\[
 n>N_1\quad\Longrightarrow\quad s_n>M/c.
\]
Since $M/c>0$, these $s_n$ are positive. If
$n>\max\{N_0,N_1\}$, we can multiply $t_n\ge c$ by the positive
number $s_n$, obtaining
\[
 s_nt_n\ge cs_n>c(M/c)=M.
\]
Thus every sufficiently late product exceeds $M$. Checking positive
thresholds suffices: exceeding any positive number also exceeds all
smaller, possibly negative, thresholds. Hence $s_nt_n\to+\infty$.

\step{(b) A limit superior of $+\infty$.}
Now assume only $\limsup s_n=+\infty$. Every tail of $(s_n)$ is
unbounded above. Indeed, a tail bounded above, together with the
finitely many preceding real terms, would bound the whole sequence
above and force its limit superior to be at most a real upper bound.

Fix a tail index $N$ and a positive threshold $M$. The unbounded
tail beyond $\max\{N,N_0\}$ contains an index $n$ for which
$s_n>M/c$. At that index $n>N$, and the same sign calculation as
in part (a) gives $s_nt_n>M$. Thus every tail of the product sequence
is unbounded above. Its tail suprema are all $+\infty$, so its
limit superior is $+\infty$.

The distinction between (a) and (b) is in the quantifiers: part (a)
puts every sufficiently late product above $M$; part (b) produces
at least one such product in every tail.

\step{(c) Recover Exercise 12.7.}
If $t_n=k>0$ for every $n$, then $\liminf t_n=k>0$ and
$s_nt_n=ks_n$. Part (b) therefore states exactly that
\[
 \limsup s_n=+\infty\quad\Longrightarrow\quad
 \limsup(ks_n)=+\infty\qquad(k>0).
\]
\answer{Both product conclusions hold; a positive constant second
factor gives Exercise~12.7.}

\exercise{12.10}{Detecting boundedness with the limit superior of absolute values}
\goal Prove both directions of
\[
 (s_n)\text{ is bounded}\quad\Longleftrightarrow\quad
 \limsup\abs{s_n}<+\infty.
\]
\idea Absolute values express upper and lower bounds on $s_n$ in one
nonnegative sequence. An eventual bound becomes a bound on all terms
after we include the finite initial segment.

\step{A bounded sequence gives a finite limit superior.}
If $(s_n)$ is bounded, there is a finite $B\ge0$ with
$\abs{s_n}\le B$ for every $n$. Thus for every $N$,
\[
 0\le\sup\{\abs{s_n}:n>N\}\le B.
\]
These tail suprema decrease as $N$ increases. The monotone sequence
theorem gives a limit in $[0,B]$, and by definition that limit is
$\limsup\abs{s_n}$. In particular it is less than $+\infty$.

\step{A finite limit superior gives an eventual bound.}
Conversely, suppose $L=\limsup\abs{s_n}<+\infty$. Since all the
absolute values are nonnegative, $L\ge0$, so $L$ is a real number.
The sequence $(\abs{s_n})$ cannot be unbounded above: if it were,
every tail would be unbounded above, since deleting finitely many
real terms cannot remove unboundedness. All its tail suprema would
then be $+\infty$, contradicting the finite value of $L$.

Equivalently, some tail has a finite supremum. Once that happens,
all subsequent tail suprema are finite and converge to $L$.
Their convergence yields an integer $N\ge1$ with
\[
 \abs{s_n}\le\sup\{\abs{s_j}:j>N\}<L+1
 \qquad(n>N).
\]

\step{Include the early terms.}
The finite set of real numbers
$\{\abs{s_1},\ldots,\abs{s_N},L+1\}$ has a finite maximum $B$.
If $n\le N$, the definition of $B$ gives $\abs{s_n}\le B$.
If $n>N$, the tail estimate gives $\abs{s_n}<L+1\le B$.
Consequently $\abs{s_n}\le B$ for every index $n$.
\answer{The two conditions are equivalent. The limit superior bounds
the tail, and a finite maximum absorbs the initial terms.}

\exercise{12.11}{The lower ratio--root inequality}
\goal For a sequence with $s_n\ne0$ at every index, prove the first
inequality of Theorem~12.2:
\[
 \liminf\abs{\frac{s_{n+1}}{s_n}}
 \ \le\ \liminf\abs{s_n}^{1/n}.
\]
\idea A lower bound on successive ratios becomes a geometric lower
bound on the terms. Taking $n$th roots removes the fixed factor.

\step{Name the two nonnegative extended limits.}
Set
\[
 A=\liminf\abs{\frac{s_{n+1}}{s_n}},\qquad
 B=\liminf\abs{s_n}^{1/n}.
\]
Every ratio is well defined and positive because no $s_n$ is zero.
The roots are positive as well. Therefore $A,B\in[0,+\infty]$.
If $A=0$, then $A\le B$ follows immediately from $B\ge0$.

\step{An arbitrary positive number below $A$ is an eventual ratio bound.}
Suppose $A>0$, and choose any real $a$ with $0<a<A$. When
$A=+\infty$, this means that $a$ can be any positive real number.
The definition of the limit inferior gives an integer $K\ge1$
such that
\[
 \abs{\frac{s_{j+1}}{s_j}}>a\qquad(j\ge K).
\]
To see this directly, some tail infimum of these ratios exceeds $a$;
each ratio in that tail is at least the infimum. Increase the starting
index to $K$ if needed.

\step{Multiply the consecutive inequalities.}
For $n>K$, cancellation of the nonzero terms gives
\[
 \abs{s_n}
 =\abs{s_K}\prod_{j=K}^{n-1}\abs{\frac{s_{j+1}}{s_j}}
 \ge \abs{s_K}a^{n-K}.
\]
The weak inequality also holds when $n=K$, since the product then
has no factors and equals $1$. Define the fixed positive constant
$C=\abs{s_K}a^{-K}$. The estimate becomes
\[
 \abs{s_n}\ge Ca^n\qquad(n\ge K).
\]
The map $x\mapsto x^{1/n}$ is increasing on nonnegative reals, so
\[
 \abs{s_n}^{1/n}\ge aC^{1/n}\qquad(n\ge K).
\]

\step{Pass to the limit inferior and then let $a$ approach $A$.}
Since $C>0$, the elementary root limit gives $C^{1/n}\to1$.
Thus $aC^{1/n}\to a$, and eventual order (Exercise~12.1) yields
$B\ge a$.

If $A$ is finite and $B<A$, then $B$ is finite and a number
$a$ strictly between $B$ and $A$ satisfies $a>0$. This contradicts
$B\ge a$. Hence $B\ge A$ in the finite case. If $A=+\infty$,
the inequality $B\ge a$ holds for every positive real $a$; a finite
$B$ would be contradicted by taking $a>B$. Therefore $B=+\infty$.
\answer{In every case the lower ratio limit inferior is at most the
root limit inferior. This proof does not rely on Corollary~12.3,
whose proof uses the inequality being established here.}

\exercise{12.12}{Arithmetic means and the limits of the original sequence}
\goal For $s_n\ge0$, put
\[
 \sigma_n=\frac{s_1+\cdots+s_n}{n}.
\]
Prove the four-term inequality in part (a), deduce part (b), and give
the counterexample requested in part (c).

\step{(a), preparation: name the bounds and check nonnegativity.}
Write $\alpha=\liminf s_n$ and $\beta=\limsup s_n$.
Because $s_n\ge0$, all averages are nonnegative, and
$\alpha,\beta,\liminf\sigma_n,\limsup\sigma_n$ belong to
$[0,+\infty]$. The inequality
$\liminf\sigma_n\le\limsup\sigma_n$ holds for every sequence
by Theorem~10.7. We prove the two outside inequalities separately.

\step{(a), upper bound: verify the estimate in the hint.}
First suppose $\beta<+\infty$. The sequence $(s_n)$ must then be
bounded above: otherwise every tail would be unbounded above, and
its limit superior would be $+\infty$. Set
\[
 U_N=\sup\{s_n:n>N\},\qquad P_N=s_1+\cdots+s_N.
\]
These are finite, with $U_N,P_N\ge0$. If $n>N$, each of the
$n-N$ remaining terms is at most $U_N$, so
\[
 \sigma_n\le\frac{P_N}{n}+\frac{n-N}{n}U_N
 \le\frac{P_N}{n}+U_N.
\]
The second inequality uses $U_N\ge0$ and $(n-N)/n\le1$.
If $n>M>N$, the inequality $P_N/n\le P_N/M$ uses $P_N\ge0$.
Taking the supremum over $n>M$ gives the hint exactly:
\[
 \sup\{\sigma_n:n>M\}\le\frac{P_N}{M}+U_N.
\]
Hold $N$ fixed and let $M\to\infty$. The left side tends to
$\limsup\sigma_n$, while $P_N/M\to0$, yielding
$\limsup\sigma_n\le U_N$. This holds for every $N$.
Letting $N\to\infty$ and using $U_N\to\beta$, we get
$\limsup\sigma_n\le\beta$. All the relevant quantities are finite
here, since a global bound on $s_n$ also bounds each average.

If instead $\beta=+\infty$, the desired upper bound is simply
$\limsup\sigma_n\le+\infty$, which always holds. No finite-tail
estimate involving an infinite supremum is required.

\step{(a), lower bound: retain a uniformly positive tail.}
If $\alpha=0$, nonnegativity already gives
$\liminf\sigma_n\ge0=\alpha$.
If $\alpha>0$, choose an arbitrary real $c$ with $0<c<\alpha$.
For $\alpha=+\infty$, take any $c>0$. There is $N$ such that
$s_j\ge c$ for all $j>N$, by the definition of the limit inferior.
For $n>N$, discard the nonnegative initial sum and estimate the tail:
\[
 \sigma_n
 =\frac{P_N+s_{N+1}+\cdots+s_n}{n}
 \ge\frac{(n-N)c}{n}
 =c\left(1-\frac{N}{n}\right).
\]
The expression on the right tends to $c$, since $N$ is fixed.
Exercise~12.1 gives $\liminf\sigma_n\ge c$. When $\alpha$ is
finite, letting $c$ range over all positive numbers below $\alpha$
forces $\liminf\sigma_n\ge\alpha$: otherwise choose $c$ between
these two quantities. When $\alpha=+\infty$, the same inequality
for every $c>0$ forces $\liminf\sigma_n=+\infty$.
Combining the three inequalities proves
\[
 \boxed{\liminf s_n\le\liminf\sigma_n\le
 \limsup\sigma_n\le\limsup s_n.}
\]

\step{(b) Convergence passes to the averages.}
Suppose $s_n\to L$. Nonnegativity means that either $L\ge0$ is
finite or $L=+\infty$; a limit of $-\infty$ is impossible.
Then $\liminf s_n=\limsup s_n=L$, and part (a) gives
\[
 L\le\liminf\sigma_n\le\limsup\sigma_n\le L.
\]
Thus both middle quantities equal $L$. Theorem~10.7, which
characterizes convergence by equality of the limit inferior and
limit superior, implies $\sigma_n\to L$. If $L=+\infty$, this
statement means that the averages also tend to $+\infty$.

\step{(c) Convergent averages need not imply convergent terms.}
Take $s_n=1+(-1)^n$. Its terms alternate $0,2,0,2,\ldots$, and
are nonnegative. The odd subsequence is constantly $0$ and the even
subsequence is constantly $2$, so $(s_n)$ has no limit.

For $n=2k$, there are $k$ zeros and $k$ twos, so
$\sigma_{2k}=2k/(2k)=1$. For $n=2k-1$, there are $k$ zeros and
$k-1$ twos, so
\[
 \sigma_{2k-1}=\frac{2(k-1)}{2k-1}
 =1-\frac{1}{2k-1}.
\]
In both cases $\abs{\sigma_n-1}\le1/n$. Given $\varepsilon>0$,
choose $N>1/\varepsilon$; then $n>N$ implies
$\abs{\sigma_n-1}\le1/n<\varepsilon$. Thus $\sigma_n\to1$.
\answer{The averages satisfy the four-term bounds and inherit every
limit of a nonnegative sequence. The converse fails for alternating
zeros and twos.}

\exercise{12.13}{Limit inferior and superior as eventual bounds}
\goal For a bounded real sequence define
\[
 \begin{split}
 A&=\{a\in\R:\{n\in\N:s_n<a\}\text{ is finite}\},\\
 B&=\{b\in\R:\{n\in\N:s_n>b\}\text{ is finite}\}.
 \end{split}
\]
Prove $\sup A=\liminf s_n$ and $\inf B=\limsup s_n$.
\idea A finite set of exceptional indices can be placed before a
single cutoff. Thus $A$ consists of eventual lower bounds and $B$
of eventual upper bounds. The boundary value itself need not belong
to the respective set.

\step{Check existence of the real suprema and infima.}
Choose $C\ge0$ with $-C\le s_n\le C$ for every $n$.
For $a=-C$, no term satisfies $s_n<a$, so $-C\in A$.
If $a>C$, every index satisfies $s_n<a$, so $a\notin A$.
Thus $A$ is nonempty and bounded above by $C$.
Similarly, $C\in B$, and every $b<-C$ lies outside $B$, since
$s_n>b$ for every $n$. Hence $B$ is nonempty and bounded below
by $-C$. Completeness gives finite real values $\sup A$ and $\inf B$.
Write
\[
 \alpha=\liminf s_n,\quad\beta=\limsup s_n,\quad
 L_N=\inf\{s_n:n>N\},\quad U_N=\sup\{s_n:n>N\}.
\]
All these numbers are finite, and $L_N\to\alpha$, $U_N\to\beta$.

\step{Every eventual lower bound is at most $\alpha$.}
If $a\in A$, its finitely many exceptional indices lie at most at
some index $N_0$; if there are none, take $N_0=1$.
Then $s_n\ge a$ whenever $n>N_0$. For $N\ge N_0$, it follows
that $L_N\ge a$, and passage to the limit gives $\alpha\ge a$.
Thus $\alpha$ is an upper bound of $A$, so $\sup A\le\alpha$.

\step{No number below $\alpha$ can be an upper bound of $A$.}
If $a<\alpha$, convergence $L_N\to\alpha$ implies $L_N>a$
for some $N$. For every $n>N$, we have $s_n\ge L_N>a$.
Consequently the indices with $s_n<a$ are contained in
$\{1,\ldots,N\}$, a finite set; hence $a\in A$.
Now fix any $d<\alpha$ and take $a=(d+\alpha)/2$.
Then $d<a<\alpha$, so $a\in A$ and $a>d$. Therefore $d$ fails
to be an upper bound of $A$. This proves that $\alpha$ is the
least upper bound: $\sup A=\alpha$.

\step{Repeat the two directions for eventual upper bounds.}
If $b\in B$, there is $N_0$ such that $s_n\le b$ whenever
$n>N_0$. Then $U_N\le b$ for all $N\ge N_0$; taking limits gives
$\beta\le b$. Thus $\beta$ is a lower bound of $B$.

If $b>\beta$, convergence $U_N\to\beta$ gives some $N$ with
$U_N<b$. For every $n>N$, we have $s_n\le U_N<b$, so the
indices satisfying $s_n>b$ form a finite set. Thus $b\in B$.
If $d>\beta$, the number $b=(\beta+d)/2$ satisfies
$\beta<b<d$ and belongs to $B$, so $d$ cannot be a lower bound
of $B$. Hence $\beta$ is its greatest lower bound:
$\inf B=\beta$.

\step{Why the boundary is not asserted to lie in the set.}
For $s_n=-1/n$, one has $\alpha=0$, but every $s_n<0$,
so $0\notin A$. For $s_n=1/n$, one has $\beta=0$, but every
$s_n>0$, so $0\notin B$. Suprema and infima need not be attained.
\answer{The least eventual upper bound is $\limsup s_n$, and the
greatest eventual lower bound is $\liminf s_n$, in the precise
supremum/infimum sense above.}

\exercise{12.14}{The growth of the geometric mean of the first integers}
\goal Calculate the limits of $(n!)^{1/n}$ and $(n!)^{1/n}/n$.
\idea Apply the ratio--root corollary to two different positive
sequences. For the second limit we also justify the elementary
sequence that defines $e$, without using Stirling's formula.

\step{(a) Apply the ratio--root corollary to the factorials.}
Let $v_n=n!$, which is positive. Cancelling the common product
$1\cdot2\cdots n$ gives
\[
 \frac{v_{n+1}}{v_n}=\frac{(n+1)!}{n!}=n+1\longrightarrow+\infty.
\]
Corollary~12.3 states that a limit of the absolute successive ratios
is also the limit of the $n$th roots, including the value $+\infty$.
Therefore
\[
 \boxed{\lim_{n\to\infty}(n!)^{1/n}=+\infty.}
\]

\step{(b), preparation: justify the sequence defining $e$.}
Set $b_n=(1+1/n)^n$. The standard constant $e$ is the limit of
this sequence, as described in the text's Section~7, Example~1(e).
For completeness we verify that this limit exists and is positive.
By the binomial theorem,
\[
 b_n=\sum_{k=0}^{n}\binom{n}{k}\frac{1}{n^k}
 =\sum_{k=0}^{n}\frac{1}{k!}
      \prod_{j=0}^{k-1}\left(1-\frac{j}{n}\right).
\]
For $k=0$ the product is empty and equals $1$.
For fixed $k\le n$, each displayed factor is nonnegative and
does not decrease when $n$ is replaced by $n+1$. Thus each of the
old summands does not decrease. The extra summand with $k=n+1$
in the expression for $b_{n+1}$ is positive, so $b_{n+1}>b_n$.

Every factor is at most $1$. For $k\ge2$, the $k-1$ factors
$2,3,\ldots,k$ in $k!$ are each at least $2$, whence
$k!\ge2^{k-1}$. It follows that
\[
 b_n\le\sum_{k=0}^{n}\frac{1}{k!}
 \le 2+\sum_{k=2}^{n}\frac{1}{2^{k-1}}<3.
\]
For $n=1$ the last sum is empty, and the same bound holds.
Since $b_1=2$, the monotone sequence theorem now gives
$b_n\to e$ for a finite number $e\ge2>0$.

\step{(b), normalization: put the denominator inside the root.}
Define $u_n=n!/n^n>0$. Its positive $n$th root is
\[
 u_n^{1/n}=\frac{(n!)^{1/n}}{n},
\]
because the $n$th power of the expression on the right is
$n!/n^n$ and the positive $n$th root is unique.
The successive ratio simplifies one factor at a time:
\[
 \begin{aligned}
 \frac{u_{n+1}}{u_n}
 &=\frac{(n+1)!}{(n+1)^{n+1}}\frac{n^n}{n!}\\
 &=\frac{(n+1)n^n}{(n+1)^{n+1}}
 =\left(\frac{n}{n+1}\right)^n
 =\frac{1}{(1+1/n)^n}
 =\frac{1}{b_n}.
 \end{aligned}
\]
Since $b_n\to e\ne0$, the reciprocal limit law gives
$u_{n+1}/u_n\to1/e$. Corollary~12.3 therefore gives
$u_n^{1/n}\to1/e$, which is exactly
\[
 \boxed{\lim_{n\to\infty}\frac{(n!)^{1/n}}{n}=\frac{1}{e}.}
\]
\answer{The unnormalized geometric mean tends to $+\infty$;
after division by $n$, it tends to $1/e$.}
